% s-352-p352-m001; original PDF 352
\Aphorism{s-352-p352-m001}{403}
Gott sieht die Welt von aussen, durch den Zeitschleier hindurch, als sinnloses heilloses hoffnungsloses Zufallsspiel der ewigen Wiederkunft. Welcher Mensch hielte das aus, ohne am Mitleiden mit der Welt zu ersticken? Wer wollte es wagen, unter dieser perspectivischen Grundbedingung Gott zu sein? Darum — giebt es keinen Gott, oder aber, er ist eines jener Ungethüme von Herzlosigkeit und \Greek{ἀταραξία}, nach dem wir die Pfeile unserer menschlich-allzumenschlichen, allzuzwergenhaften Worte, Wünsche und Werthschätzungen umsonst abschiessen, ein epikureischer Zwischengott, der sich von der Lenkung der Dinge zurückgezogen hat, weil für ihn die Welt bereits unendlich oft vergangen, also starr und incorrigibel ist.
\par

